\section{The full dyadic image}
\label{sec:full-dyadic-image}

The normalized image and its cyclotomic action determine the
whole group acting on the free pro-two fundamental group.
Write
\[
 Q_2=\operatorname{im}\bigl(G\longrightarrow
                 \operatorname{Aut}(\widehat F_2^{(2)})\bigr).
\]
The symbol \(Q_2\) denotes a group, not the field \(\mathbf Q_2\).

\begin{theorem}\label{thm:full-dyadic-free-product}
The full dyadic image is the arithmetic image, and
\[
                  Q_2\cong C_2*_{\mathrm{pro}\text{-}2}\mathbf Z_2.
\]
One may take complex conjugation as the generator of \(C_2\),
and any arithmetic element with dyadic multiplier \(-3\) as
the generator of the second factor.
Under this presentation the multiplier is the abelianization
map
\[
 Q_2\longrightarrow\mathbf Z_2^\times
        =\{\pm1\}\times(1+4\mathbf Z_2),
 \qquad c\longmapsto-1,\qquad h\longmapsto-3.
\]
Its kernel is the normalized image \(P=M/R_M\).
\end{theorem}

\subsection{Two arithmetic generators}

\begin{proof}[Generation and the cyclotomic kernel]
Every dyadic action is the identity on the mod-two Frattini
quotient of the free group. The automorphisms with this property
form a pro-two group, by the finite characteristic filtration
argument used for the prime word actions. Thus \(Q_2\) is pro-two.
Uniform arithmetic saturation gives equality of its full and
arithmetic images.

We first check that its cyclotomic kernel is exactly \(P\).
The tangential arithmetic action is unramified at every odd
prime \cite[\S8.5, p.~151]{HainMatsumoto2003}.
For an odd prime \(q\), inertia at \(q\) therefore acts trivially
on \(\widehat F_2^{(2)}\), whereas its cyclotomic image is the
whole \(q\)-local factor \(\mathbf Z_q^\times\) and is trivial
at every other prime. The closed subgroup generated by these
inertia groups has trivial dyadic action. Its full multiplier
image contains the finite-support product of the odd-local
unit groups, and compactness makes that image the whole product.
Consequently an arithmetic element with dyadic multiplier one
can have all its other multiplier coordinates cancelled without
changing its dyadic action. Its resulting lift belongs to
\(M\cap H\). This proves the kernel assertion.

Let \(c\) be complex conjugation, using the positive tangential
base point. Then \(c^2=1\), \(f_c=1\), and \(\rho_2(c)=0\).
Choose \(h\in H\) with \(\lambda_2(h)=-3\) and set
\(a_3=[c,h]\). The cocycle law gives
\[
 \rho_{2,2}(a_3)=-2\rho_{2,2}(h),\qquad
 \mu(a_3)=-\rho_{2,2}(h).
\]
The parity theorem says that \(\rho_{2,2}(h)\) is odd, since
\(-3\equiv5\pmod8\). Hence the cubic leading coefficient of
\(a_3\) is odd. Its reduction is the prescribed weight-three
free Lie generator.

Apply the uniform generator construction with this seed and
this same \(h\). Every constructed generator belongs to the
closed subgroup generated by \(c,h\): the construction uses
conjugation by \(h\), powers, and corrections made from earlier
generators. Its leading terms give the complete integral basis
in every weight. Successive correction and compactness therefore
put all of \(P\) in that closed subgroup.
The multipliers \(-1,-3\) generate \(\mathbf Z_2^\times\), so
\(c,h\) generate \(Q_2\).

For these generators the construction also gives
\[
 [a_{3+2j}]=(h-1)^j[a_3]\quad\hbox{in }P/\Phi(P).
\]
Indeed the conjugation difference in odd weight \(n\) differs
from \([h,a_n]\) by a power with even dyadic exponent, and
the further corrections are squares. All disappear in the
Frattini quotient. The arithmetic leading-basis theorem and
the normalized freeness theorem apply to this choice of seed.
\end{proof}

\subsection{A free kernel with a convergent basis}

We give the pro-two group calculation needed to detect relations.
Let
\[
 F=C_2*_{\mathrm{pro}\text{-}2}\mathbf Z_2
   =\langle s,t\mid s^2=1\rangle_{\mathrm{pro}\text{-}2},
 \qquad
 \pi:F\longrightarrow C_2\times\mathbf Z_2
\]
be the map sending \(s\) and \(t\) to the two factors, and put
\(K_F=\ker\pi\).

\begin{lemma}\label{lem:free-product-cyclotomic-kernel}
The group \(K_F\) is free pro-two on the convergent sequence
\[
 b_0=[s,t],\qquad b_{j+1}=[t,b_j]\quad(j\ge0).
\]
\end{lemma}

\begin{proof}
The kernel of \(F\to C_2\) is free pro-two on
\(t,k=sts\). This follows directly from universal properties:
if \(A\) is free pro-two on \(t,k\), the involution interchanging
these generators defines \(A\rtimes C_2\); the maps in both
directions given by \(s,t,k=sts\) are inverse.
Replacing \(k\) by \(b_0=kt^{-1}\) is a free change of basis.
Thus \(K_F\) is the kernel of the map from the free pro-two
group on \(t,b_0\) to \(\mathbf Z_2\) that sends \(t\) to one
and \(b_0\) to zero.

Let \(\mathcal P\) be free pro-two on the sequence
\(z_0,z_1,\ldots\) converging to one, concretely the inverse
limit of the free pro-two groups on its finite initial segments
with retractions killing the last generator. Define
\[
                   \alpha(z_j)=z_{j+1}z_j.
\]
On every finite initial-segment quotient this is an invertible
triangular change of free basis: solve for the generators
starting from the last one. These automorphisms and their
inverses commute with the retractions. They therefore define
an automorphism of \(\mathcal P\).
Its action on each finite-rank Frattini quotient is unipotent.
The kernel of the action of the automorphism group of a
finite-rank free pro-two group on its Frattini quotient is
pro-two. Hence the closed cyclic group generated by each
finite-rank image of \(\alpha\) is pro-two. Passing to the
inverse limit shows that powers of \(\alpha\) define a
continuous \(\mathbf Z_2\)-action on \(\mathcal P\).

Form \(\mathcal P\rtimes_\alpha\mathbf Z_2\), writing \(t\)
for the second-factor generator. It is generated by \(t,z_0\),
because \([t,z_j]=z_{j+1}\). The free pro-two group on
\(t,b_0\) maps onto it by \(b_0\mapsto z_0\).
Conversely, in that free group the recursively defined
\(b_j\) converge to one: in every finite two-group quotient,
iterated commutators eventually vanish by nilpotence.
The universal property gives a continuous map
\(\mathcal P\to A\), \(z_j\mapsto b_j\).
The recursion makes it equivariant for conjugation by \(t\),
and therefore for the full \(\mathbf Z_2\)-action by continuity.
It extends to a map of semidirect products into \(A\).
The two maps fix the generators \(t,b_0\), respectively
\(t,z_0\), and so are inverse.
The kernel of the exponent map is thus precisely
\(\mathcal P\), proving the lemma.
\end{proof}

\subsection{There are no further relations}

\begin{proof}[Completion of Theorem~\ref{thm:full-dyadic-free-product}]
Generation gives a surjection
\[
                   F\longrightarrow Q_2,
                   \qquad s\longmapsto c,\quad t\longmapsto h.
\]
It commutes with the cyclotomic quotient maps, so restricts
to a surjection \(K_F\to P\).
The preceding lemma and the exact Frattini recursion identify
its map on Frattini quotients with the identity on the
products of copies of \(\mathbf F_2\):
\[
             [b_j]\longmapsto[a_{3+2j}].
\]

We verify carefully that this forces injectivity here.
The surjection \(K_F\to P\) has a continuous section, since
\(P\) is free on a convergent basis. To construct one, choose
a descending countable neighborhood basis of open subgroups
of \(K_F\). Their images are open in \(P\). Late basis
elements of \(P\) lie in each such image. They can therefore
be lifted so that the chosen lifts converge to one; choose
arbitrary lifts for the finitely many initial exceptions.
The free universal property produces a section.
Thus \(K_F=J\rtimes P\), where \(J\) is the kernel.

If \(J\ne1\), it has a nonzero continuous character to
\(\mathbf F_2\), obtained from a nontrivial finite two-group
quotient. The orbit of this character under \(P\) is finite:
continuity and compactness of \(J\) give an open stabilizer.
Its span is a nonzero finite-dimensional \(\mathbf F_2\)-module
for a finite two-group. Such a module has a nonzero invariant
vector. For example, induct on the order of the group, first
taking the nonzero kernel of \(z-1\) for a central element
\(z\) of order two and then taking invariants under the
quotient group.
The resulting invariant character of \(J\) extends to
\(J\rtimes P\) by setting it zero on \(P\).
This contradicts the isomorphism of Frattini quotients, which
says that every character of \(K_F\) is pulled back from \(P\)
and hence vanishes on \(J\). Therefore \(J=1\).

The map \(F\to Q_2\) is injective on the kernel of its
cyclotomic quotient and is an isomorphism on that quotient.
It is consequently an isomorphism. Finally,
\(F^{\mathrm{ab}}=C_2\times\mathbf Z_2\), by the universal
property of the free product. This proves the abelianization
assertion as well.
\end{proof}

This theorem describes one fundamental-group realization.
It does not identify the full profinite group \(G\) with
its arithmetic subgroup or determine the residual quotient
supported on \(R_D\).
